\documentclass[11pt,a4paper]{article}

\usepackage[utf8]{inputenc}
\usepackage[T1]{fontenc}
\usepackage[spanish]{babel}
\usepackage[margin=2cm, bottom=6.2cm, footskip=0.6cm]{geometry}
\usepackage{graphicx}
\usepackage{tabularx}
\usepackage{array}
\graphicspath{{../../01_Propuesta/Logos/}{./}}
\usepackage{fancyhdr}
\usepackage{xcolor}
\usepackage[hidelinks]{hyperref}
\usepackage{enumitem}
\setlist{nosep}
\setlength{\parindent}{0pt}
\setlength{\parskip}{4pt}

\definecolor{verdeosc}{RGB}{26,107,60}
\definecolor{verdemid}{RGB}{46,139,87}

\pagestyle{fancy}
\fancyhf{}
\fancyfoot[C]{%
  \tiny
  \hyphenpenalty=10000\exhyphenpenalty=10000
  \noindent
  \begin{minipage}{\textwidth}
  \begin{minipage}[t]{4.4cm}
  {\color{verdeosc}\textbf{ORGANIZA}}\\[3pt]
  Grupo de Investigación IN SILICO\\[6pt]
  {\color{verdeosc}\textbf{CO-ORGANIZADORES}}\\[3pt]
  Departamento de Biología y Química\\
  Departamento de Ingeniería Agroindustrial\\
  Departamento de Ingeniería Agrícola\\
  Departamento de Física\\
  Universidad de Sucre
  \end{minipage}%
  \hfill
  \begin{minipage}[t]{3.7cm}
  {\color{verdeosc}\textbf{COORDINACIÓN}}\\[3pt]
  \textbf{Aldo F. Combariza}\\
  Coordinación General\\
  \textbf{María Ximena Díaz}\\
  Coordinación Académica\\
  \textbf{Sebastián Vargas}\\
  Coordinación Logística\\
  \textbf{Selena Arias Avila}\\
  Coordinación Administrativa
  \end{minipage}%
  \hfill
  \begin{minipage}[t]{3.7cm}
  {\color{verdeosc}\textbf{CONTACTO}}\\[3pt]
  \textbf{SICVEC 2026}\\
  19--20 de octubre de 2026\\
  C.C. Guacarí, Sincelejo\\
  Modalidad híbrida\\
  \href{mailto:insilico@unisucre.edu.co}{insilico@unisucre.edu.co}
  \end{minipage}%
  \end{minipage}%
}
\renewcommand{\headrulewidth}{0pt}

\begin{document}

\noindent
\includegraphics[height=1.6cm]{logo_sicvec2026.png}%
\hspace{10pt}{\color{verdeosc}\vrule width 1.2pt height 1.5cm}\hspace{10pt}%
\includegraphics[height=1.28cm]{Logo_insilico.jpeg}%
\hspace{10pt}{\color{verdeosc}\vrule width 1.2pt height 1.5cm}\hspace{10pt}%
\includegraphics[height=1.3cm]{Logo_Bio.png}\hspace{6pt}%
\includegraphics[height=1.3cm]{logo_agroin.png}\hspace{6pt}%
\includegraphics[height=1.3cm]{logo_ingagrocola.png}\hspace{6pt}%
\includegraphics[height=1.3cm]{logo_fisica_lifi.png}%
\hspace{10pt}{\color{verdeosc}\vrule width 1.2pt height 1.5cm}\hspace{10pt}%
\includegraphics[height=1.3cm]{logo_unisucre.png}

\vspace{0.4cm}

\noindent
\textbf{Sincelejo, 7 de octubre de 2026}

\vspace{0.15cm}

\noindent
\textbf{A los Jefes de los Departamentos co-organizadores}\\
Departamento de Biología y Química · Departamento de Ingeniería Agroindustrial\\
Departamento de Ingeniería Agrícola · Departamento de Física\\
Universidad de Sucre

\vspace{0.15cm}

\noindent
\textbf{\large Asunto: Información final y solicitud de difusión — SICVEC 2026}

\vspace{0.15cm}

Estimados colegas:

\vspace{0.1cm}

A pocos días del simposio, les compartimos la información final de SICVEC 2026 y les pedimos un favor: que la difundan por los canales de su departamento, entre sus estudiantes, docentes y redes académicas e investigativas.

\vspace{0.1cm}

\noindent
\textbf{INFORMACIÓN FINAL DEL EVENTO}

\begin{tabularx}{\linewidth}{@{}l X@{}}
\textbf{Evento:} & Simposio Internacional de Ciencia Verde y Economía Circular (SICVEC 2026) \\
\textbf{Fechas:} & 19 y 20 de octubre de 2026 (lunes y martes) \\
\textbf{Lugar:} & Centro Comercial Guacarí, Sincelejo, Sucre \\
\textbf{Modalidad:} & Híbrida — presencial por categoría / virtual gratuita (externos e internacionales) \\
\textbf{Cierre de resúmenes:} & 12 de octubre de 2026, 23:59 (hora Colombia) \\
\textbf{Evaluación de pares:} & hasta el 15 de octubre de 2026, 23:59 \\
\textbf{Decisión y notificación:} & 15 de octubre de 2026 \\
\textbf{Límite de pago (presencial):} & 18 de octubre de 2026, 23:59 \\
\textbf{Sitio web (en español):} & \href{https://sicvec2026.eu.pythonanywhere.com/?lang=es}{sicvec2026.eu.pythonanywhere.com/?lang=es} \\
\textbf{Inscripción virtual (gratuita):} & \href{https://sicvec2026.eu.pythonanywhere.com/inscripcion?lang=es}{sicvec2026.eu.pythonanywhere.com/inscripcion?lang=es} \\
\end{tabularx}

\vspace{0.1cm}

Adjuntamos el afiche oficial del evento para que lo usen en la difusión, junto con esta carta.

\vspace{0.1cm}

Muchas gracias por su apoyo como departamentos co-organizadores, y por ayudarnos a que SICVEC 2026 llegue a la mayor cantidad posible de estudiantes y colegas.

\vspace{0.15cm}

Atentamente,

\vspace{0.15cm}

\textbf{Dr. Aldo F. Combariza}\\
Coordinador General, SICVEC 2026\\
Grupo de Investigación IN SILICO · Departamento de Biología y Química\\
Facultad de Educación y Ciencias · Universidad de Sucre\\
\href{mailto:aldo.combariza@unisucre.edu.co}{aldo.combariza@unisucre.edu.co}

\end{document}
